\section{Generic chains and orders of sections}
\label{sec:chains}

The canonical-index argument will produce curves of uniformly bounded
degree divided by their marked multiplicity on finite product covers.
We need to replace those curves by the fibres of a rational map and
retain quantitative degree bounds under projection and finite quotient.
This section carries out that step. Its order propagation uses the
parameter-differentiation method of Ein--K\"uchle--Lazarsfeld
\cite[Proposition~2.3]{EinKuchleLazarsfeld1995}, applied here along
generic chains. For the quotient we use the stabilization strategy in
Campana's chain-equivalence construction
\cite[Theorem~1.1 and its proof sketch]{Campana2004}, with the
generic-point formulation and quantitative estimates required here.
We prove the projection and finite-quotient compatibilities as well.

We work in characteristic zero and first specify the generic-point
conventions used in the arguments.  All finite collections of varieties,
line bundles, morphisms, and sections over $\C$ descend to countable
subfields.  After enlarging such a field and taking its algebraic closure,
we may take successive geometric generic points inside a sufficiently
large algebraically closed extension.  A point chosen generically at a
particular stage is generic over the field of all the data already chosen.
We always specify a smaller field when taking the closure of that point.

Conditions on jet-kernel dimensions involve only countably many pairs of
integers: a degree and an order.  For each pair, generic base change and
semicontinuity on smooth opens give the corresponding generic condition.
Thus their simultaneous validity at a geometric generic point is
consistent with the usual very-general-point convention over $\C$.
Countably many exceptional closed sets can be included in the initial
field of definition.  Conversely, finite data at a geometric generic
point, including a section or a marked curve, can be realized over $\C$
while preserving its genericity over that field.

We also use the following spreading convention.  If new finite-type data
are chosen on a geometric generic member of a family, they are defined
over a field finitely generated over its function field.  Spreading that
field gives a \emph{dominant} extension of the parameter space, and the
new geometric generic fibre realizes the chosen data.  Resolutions and
graphs can be spread in this fashion.  Section spaces and their jet kernels
commute with algebraically closed field extension.  No assertion about
constancy of all-degree section invariants on special fibres is intended.

\subsection{The rational quotient generated by marked families}

Fix an integral projective variety $V$. A \emph{marked covering family}
on $V$ consists of an integral parameter
space $S$, a dominant marked-point map $u:S\to V$, and a family of
subvarieties $C_b\subset V$ containing $u(b)$.  After restricting to a
dense open of $S$, we require its geometric generic member to be integral
of positive dimension.  We use the reduced incidence in $S\times V$;
its endpoint evaluation dominates $V$ because it contains the graph of
$u$.  All generic choices below are made on these opens.

Starting from a generic point $a$ of $V$, a \emph{generic movement}
chooses one of the marked families, a generic point $b$ of a geometric
component of $u^{-1}(a)$ over the field of the existing path, and then a
generic endpoint on $C_b$ over that field together with $b$.

\begin{lemma}[Generic chains]\label{lem:generic-chain}
Let a nonempty finite list of marked covering families be given on an
integral projective variety $V$ in characteristic zero.  There is a
dominant rational map
\[
 \xi:V\dashrightarrow M
\]
with geometrically integral generic fibre, whose fibre closures
$H_p\ni p$ have a constant positive dimension $e$ for general $p$.
Starting at a generic $p$, a path of at most $e$ generic movements reaches
a point generic in $H_p$ over the algebraic closure of the field of $p$.
Every generic movement preserves the value of $\xi$.  In particular, a
generic family member $C_b$ is contained in $H_{u(b)}$.

The assignment is equivariant under any finite automorphism group
preserving the family list.
\end{lemma}

\begin{proof}
Choose an algebraically closed field of definition $k_*$ and a point $p$ generic in $V$ over $k_*$.  For a tuple
$a$ of points write $K_a=\overline{k_*(a)}$ inside one sufficiently large
algebraically closed extension.  Fields for successive choices include
the entire preceding path.

Every point encountered in a generic path is individually generic in
its ambient space over $k_*$.  Indeed, over a generic mark of $V$, all
components of the geometric generic fibre of a dominant map $S\to V$
have dimension $\dim S-\dim V$.  A parameter generic in such a component
is therefore generic in $S$.  The endpoint is then generic in the
incidence with that parameter, and hence in $V$.

We will compare paths after enlarging $K_p$.  The closure of a finite
path tuple over the algebraically closed field $K_p$ is geometrically
integral.  A generic point of its base change realizes the same path
after an algebraically closed field extension, with the same projected
loci.  For a fixed sequence of families, a prefix of length $j$ has total
transcendence degree
\[
 \sum_{i=1}^j
 \bigl(\dim S_i-\dim V+\dim C_{b_i}\bigr)
\]
over its initial-point field: these are precisely the increments from
the parameter and endpoint choices.  Any tuple satisfying those
incidences has at most this transcendence degree.  A path generic after
an algebraically closed field extension has the displayed degree there,
and cannot have a smaller degree before the extension.  Thus all prefix
degrees agree with their upper bounds, proving the converse comparison
as well.  This observation also applies when a path is begun at a later
endpoint and earlier path data are subsequently adjoined to its ground
field.

Let $d_j$ be the largest dimension of the closure \emph{over $K_p$} of
an endpoint after $j$ generic movements; set $d_0=0$.  These loci can be
obtained from components of the finite-step path spaces and their
projections.  Appending a step makes the new endpoint locus contain the
old one.  To see this, the new locus contains, after extension to the
field of the preceding path and the new parameter, the whole integral
member of that last step.  It therefore contains the previous endpoint,
and hence its closure over $K_p$.  Consequently $d_j$ is nondecreasing,
$d_1>0$, and the first equality has the form
\[
 d_h=d_{h+1}=e,\qquad 1\le h\le e\le\dim V.
\]
Choose an endpoint locus $H_p$ of dimension $e$ at length $h$ and a
corresponding generic endpoint $y$.

Every possible generic next movement, made over the preceding path,
has its member and endpoint in $H_p$: its endpoint locus contains
$H_p$, and the equality of dimensions forces equality.  By the
field-extension observation this is also true for all generic choices
over $K_{p,y}$, without retaining the rest of the path.  The new endpoint
is again generic in $H_p$ over $K_p$.  For each family and each geometric component of its parameter fibre
over the generic point of $H_p$, choose the new parameter independently
of the previous path. The stabilized dimension calculation says that
the endpoint image of this generic incidence is contained in $H_p$.
This is an identity on the incidence over $H_p$, so it persists after
each subsequent independent generic base change. Induction therefore
keeps all later paths in $H_p$, with endpoints generic there over $K_p$.
We use fresh generic realizations in this argument; a previously chosen
endpoint need not remain generic after arbitrary old data are adjoined.

Now form endpoint loci of paths intrinsically over $K_y$, and realize
their generic choices also over $K_{p,y}$.  They are contained in $H_p$.
Because $y$ is generic in $V$ over $k_*$, their maximum dimension at
length $h$ is the same number $e$.  Every maximizing locus therefore
becomes $H_p$ after extension.  This proves uniqueness of the maximizing
locus at $y$.  Transporting the statement between generic points of $V$
proves uniqueness at a generic starting point and shows
\begin{equation}\label{eq:chain-leaf-equality}
 H_y=H_p
 \quad\text{when }y\text{ is generic in }H_p\text{ over }K_p.
\end{equation}

Uniqueness gives descent to $k_*(p)$.  More explicitly, choose a
projective embedding over $k_*$ and take the equation subspaces of the
saturated homogeneous ideal of $H_p$ through a degree generating that
ideal.  Every algebraic conjugate is again the unique maximizing locus.
Its Grassmann coordinates are therefore invariant, and belong to
$k_*(p)$ in characteristic zero.  They define a rational map $\xi$ to
the closure of their image, naming $H_p$.  We may replace the image by
its normalization.  Equality of names means equality of subvarieties.

For a next step from the point $y$ just considered, both endpoints are
generic in $H_p$ over $K_p$, so \eqref{eq:chain-leaf-equality} gives
equality of their names.  Realizing each choice generically after field
extension shows that this equality holds at a generic parameter and
endpoint of \emph{each} family.  It is therefore a rational identity on
its incidence.  It follows that every generic movement from a generic
point preserves $\xi$, including movements that are not part of the
original maximizing path.

For completeness, fix a generic name $m=\xi(p)$ and work over $K_m$.
A generic point of any component of the geometric fibre of $\xi$ is
generic in $V$ over $k_*$.  It belongs to the subvariety bearing its own
name, so each such component is contained in the integral subvariety
named $m$.  Conversely, a generic point of that subvariety, chosen also
generically over $K_p$, has name $m$ by
\eqref{eq:chain-leaf-equality}.  Thus the closures of these fibres equal
that integral subvariety.  Characteristic zero gives geometric
reducedness as well.  The graph closure supplies a rational family with
geometrically integral generic fibre; geometric integrality and constant
dimension spread to an open of the base, by generic flatness and the
constructibility of geometric fibre properties
\cite[\S\S9,12]{EGAIV3}.  This proves the fibre assertions.  Finally, a
finite automorphism preserving the family list takes a maximizing locus
to a maximizing locus.  Uniqueness proves equivariance.  The induced rational action on the
name space can, if desired, be made regular by taking the closure of its
graph in the product of the finitely many translates.  This leaves the
generic leaves unchanged.
\end{proof}

\subsection{Projection and finite-quotient compatibility}

\begin{corollary}\label{cor:chain-functoriality}
Suppose a dominant rational map between two ambient varieties sends
generic movements of an upper family list to generic movements of a
lower family list, or to stationary steps.  Then the image of a generic
upper leaf is contained in the corresponding lower leaf.

If a finite group preserves a family list on $V$ and
$\theta:V\to V'$ is its quotient, the subvarieties
$\theta(H_p)$ are the closures of the geometrically integral generic
fibres of a rational fibration on $V'$.
\end{corollary}

\begin{proof}
The lower rational name is unchanged along every projected step.
Following a path that fills the upper leaf proves the first assertion.
It is enough here that equality of the lower names holds at the generic
point of each projected incidence: that rational identity persists when
the step is realized over the additional upper path data.

For the second assertion, let $p'$ be generic in $V'$ and lift it to
$p\in V$.  All lifts are conjugate, and equivariance gives the same
image
\[
 W_{p'}=\theta(H_p)
\]
for every lift.  Because $p$ is algebraic over $p'$, their algebraic closures in the fixed
overfield agree: $K_p=K_{p'}$. The image equation subspaces descend
to $k_*(p')$ by conjugation invariance, just as in the proof of
Lemma~\ref{lem:generic-chain}.  They give a rational name for $W_{p'}$.
A point generic in $W_{p'}$ over $K_{p'}$ lifts to a point generic in
$H_p$ over $K_p$, and hence has the same name as $p'$.  Conversely every
ambient generic point belongs to the subvariety with its own name.
The preceding fibre-of-the-name argument proves geometric integrality
and identifies its fibre closure with $W_{p'}$.
\end{proof}

\subsection{From marked-member bounds to leaf degree}

The quotient and its compatibilities depend only on the marked families.
We now put a polarization on the ambient variety. The estimate below
bounds the degree of a leaf by propagating section orders along a path
that fills it. We first define the section-order invariant used in this
argument and later in the scalar estimates.

\begin{definition}\label{def:orders}
Let $V$ be an integral projective variety of positive dimension, and let
$P$ be a big, nef, semiample rational Cartier divisor on $V$.  For a smooth
projective birational resolution $\pi:\widetilde V\to V$, set
\[
 \gamma(P;V)=
 \sup_{\substack{k>0\text{ sufficiently divisible}\\
                    0\ne s\in H^0(\widetilde V,k\pi^*P)}}
            \frac{\ord_x(s)}{k},
 \qquad
 \rho(P;V)=(P^{\dim V})^{1/\dim V},
\]
where $x$ is very general in the smooth isomorphism locus, identified
with its lift to $\widetilde V$.  If a specified
smooth point $x$ is used instead, write $\gamma(P;V,x)$, with the
resolution chosen to be an isomorphism near $x$.  For a subvariety $A$,
the notation $\gamma(P;A)$ means the same invariant for the restriction
of $P$, pulled back to a resolution of $A$.
\end{definition}

These definitions do not depend on the chosen smooth resolution: proper
birational pushforward of the structure sheaf between normal varieties
and the projection formula preserve the pulled-back section spaces.
Common resolutions then compare any two choices.  They also allow any
additional fixed divisibility requirement on $k$, since taking a power of
a section preserves its normalized order.  The invariant is finite:
on a smooth model, restrict a given nonzero section to a general
complete-intersection curve through the mark, chosen not to lie in its
zero divisor.  Its order is at most the degree of the restricted line
bundle, which is a fixed constant times $k$.  The invariant is positively
homogeneous under rational scaling of $P$ and is unchanged by rational
linear equivalence at a very general mark.

We will repeatedly use
\begin{equation}\label{eq:orders-volume}
 \gamma(P;V)\ge \rho(P;V).
\end{equation}
Indeed, put $n=\dim V$ and choose $q>0$ such that $q\pi^*P$ is Cartier
and basepoint-free on $\widetilde V$. Its morphism $f$ to its projective
image is generically finite by bigness. The Hilbert polynomial of
$f_*\OO_{\widetilde V}$, twisted by the hyperplane bundle of the image,
gives
\[
 h^0(\widetilde V,qt\pi^*P)
   =\frac{P^n}{n!}(qt)^n+O(t^{n-1}).
\]
The leading coefficient is the generic rank of
$f_*\OO_{\widetilde V}$ times the degree of the image, divided by $n!$;
it equals $(q\pi^*P)^n/n!$ by the projection formula.
Vanishing to order at least $l$ at a smooth point imposes at most
$\binom{n+l-1}{n}$ conditions.  Taking $l/k$ below, and then arbitrarily
close to, $(P^n)^{1/n}$ proves \eqref{eq:orders-volume}.

\begin{lemma}[Orders and degrees of chain leaves]\label{lem:chain-order}
Let the marked families and their rational quotient be as in
Lemma~\ref{lem:generic-chain}, and let $e$ be the dimension of a generic
leaf. Suppose that $P$ is a big, nef, semiample rational Cartier
divisor on $V$.  Assume that, for every family, its geometric generic
marked member satisfies one of the following bounds, with a common
number $B\ge0$:
\begin{enumerate}[label=\textup{(\roman*)}]
 \item the mark is smooth on $C_b$ and
       $\gamma(P;C_b,u(b))\le B$;
 \item $C_b$ is a curve and
       $P\cdot C_b/\mult_{u(b)}C_b\le B$.
\end{enumerate}
Then on the geometric generic leaf,
\begin{equation}\label{eq:chain-order-degree}
 \gamma(P;H_p)\le eB,
 \qquad
 P^e\cdot H_p\le(eB)^e.
\end{equation}
\end{lemma}

\begin{proof}
Choose an algebraically closed field of definition $k_*$ for the data.
Let $m$ be a generic point of the quotient base and put
$K_m=\overline{k_*(m)}$ inside the fixed algebraically closed overfield.
The restriction of $P$ to the geometric
generic leaf over $K_m$ is big: the leaf sweeps $V$, so it meets the generically
finite locus of the morphism defined by a basepoint-free multiple of
$P$.  Take a smooth projective resolution $J$ of this leaf over $K_m$.
The generic marked families and paths lift to $J$ wherever the resolution
is an isomorphism.  A point generic in the leaf over $K_m$ is generic in
$V$ over $k_*$.  For a generic mark $p$ on this leaf, put $K_p=\overline{k_*(p)}$.
Its name $m$ belongs to $k_*(p)$. Lemma~\ref{lem:generic-chain} supplies
a path of length $h\le e$ reaching a point generic in the whole leaf over
$K_p$.  Every
intermediate mark and endpoint is generic in $J$ over $K_m$: it is
generic in $V$ and has name $m$, and consequently has transcendence
degree $e$ over $K_m$.  In particular all their evaluation maps dominate
$J$.

Suppose that, in a divisible degree $k$, a nonzero section of $kP$ has
order at least an integer $m_0>hkB$ at the variable generic mark $p$ of
$J$.  In a fixed $K_m$-basis of $H^0(J,kP)$, the kernel of the
generic jet-evaluation matrix lets us choose such a section
$\sigma(p)$ with coefficients rational in $p$.  Follow a path of length
$h$ that fills the leaf generically over $K_p$.

Assume inductively that the same section has multiplicity at least
$m_j$ at the current mark.  Let $T$ be the closure \emph{over $K_m$}
of the path tuple through the parameter $b$ chosen for the next step;
the original point $p$ varies in this closure.  Let
$I\subset T\times J$ be the incidence obtained by adjoining the new
endpoint.  Both the graph of the current mark and $I$ dominate $J$.
Over the generic point of $T$, the generic fibre of $I$ is the chosen
member, lifted to $J$.

Differentiate in parameter directions, keeping the variable on $J$
fixed.  Concretely, use derivations of $K_m(T)$ over $K_m$ on the
coefficients of $\sigma$ in the fixed section basis.  A line-bundle
frame from $J$ is regular and nonvanishing at the generic point of $J$,
so it is unaffected by these derivations.  At the graph of the moving
mark, the multiplicity assumption is membership in the $m_j$th power
of the graph ideal, after localizing over the generic parameter.
Parameter differentiation drops its powers by at most one.  This can
also be tested after field extension to make the mark rational, using
faithful flatness.

Put
\[
 r_j=m_j-\lfloor kB\rfloor.
\]
Every iterated derivative of length less than $r_j$ still has order
strictly greater than $kB$ at the moving mark.  Each such derivative
vanishes identically on the chosen member.  In case (i), a nonzero
restriction, pulled back to its smooth resolution, would violate the
assumed order bound at the smooth mark.  The resolutions used in this
comparison agree near the mark and can otherwise be compared on a
common resolution.

In case (ii), a section with order at least $l$ at the ambient mark
restricts to an element of the $l$th power of the maximal ideal of the
local curve ring $A$.  On its normalized branches its order is at least
$l\,v_i(\mathfrak m_A)$.  The identity
\[
 \mult A=\sum_i v_i(\mathfrak m_A)
\]
then shows that a nonzero restriction has degree at least
$l\mult A$.  One way to verify this identity is to use the finite
semilocal normalization $\bar A$.  Its quotient by
$\mathfrak m_A^t\bar A$ has length
$t\sum_i v_i(\mathfrak m_A)$ over the algebraically closed residue
field.  The conductor gives a fixed $c$ such that
\[
 \mathfrak m_A^{t+c}\bar A
 \subset \mathfrak m_A^t A
 \subset \mathfrak m_A^t\bar A.
\]
Comparing leading terms in the Hilbert--Samuel lengths gives the
identity.  Here $l>kB$, whereas the degree of a nonzero restricted
section is $kP\cdot C_b\le kB\mult A$, a contradiction.

Thus all parameter derivatives of length less than $r_j$ vanish at the
generic point of $I$.  We explain why this forces order at least $r_j$
along that incidence.  Let
\[
 R=\OO_{T\times J,\eta_I},\qquad F=K_m(J).
\]
Since $I$ dominates both factors, this ring localizes over their generic
points.  It is regular, using a smooth open of $T$.  The conormal map
\[
 \mathfrak m_R/\mathfrak m_R^2
 \longrightarrow \Omega_{R/F}\otimes\kappa(I)
\]
is injective.  Indeed the residue extension is separable in
characteristic zero; in the cotangent sequence the dimensions of the
ambient differentials and residue differentials differ by $\dim R$,
which is the conormal dimension.  Equivalently, this is the generic
smoothness computation for the incidence over $J$.
Parameter derivations span the dual of this conormal after reduction,
since $R$ is a localization of
$K_m(T)\otimes_{K_m}F$.  Regularity identifies the associated graded
ring with the symmetric algebra of the conormal.  A nonzero initial
form of degree $a<r_j$ is detected by a length-$a$ iterated parameter
derivative, by the symmetric-power pairing and the invertibility of
$a!$.  The derivative vanishings therefore imply that our section, in
the chosen frame, lies in $\mathfrak m_R^{r_j}$.

Pass to the generic parameter field and then extend it so that the
independently generic endpoint is rational. The map from the ambient
product to this fibre induces a local homomorphism from $R$ to the local
ring of $J$ at that endpoint. A denominator outside the incidence prime
has nonzero value at its generic point and stays a unit after this
extension. The incidence ideal maps into the endpoint maximal ideal;
its $r_j$th power therefore maps into that maximal ideal's $r_j$th power.
The same section consequently has multiplicity at least
\[
 m_{j+1}=r_j=m_j-\lfloor kB\rfloor
\]
at the next endpoint.  After $h$ steps it still has positive
multiplicity, since $m_0>hkB$.  But the endpoint is generic in $J$ over
$K_p$, while $\sigma(p)$ is a nonzero section defined over $K_p$.
This is impossible.  Hence $\gamma(P;H_p)\le hB\le eB$.
Applying \eqref{eq:orders-volume} on the leaf gives the second inequality
of \eqref{eq:chain-order-degree}.
\end{proof}

